\documentclass[../main.tex]{subfiles}

\begin{document}
\section{Conclusion and Future Work} \label{conclusion}

In this paper, we implemented BERT for the financial domain by further pre-training it on a financial corpus and fine-tuning it for sentiment analysis (FinBERT). This work is the first application of BERT for finance to the best of our knowledge and one of the few that experimented with further pre-training on a domain-specific corpus. On both of the datasets we used, we achieved state-of-the-art results by a significant margin. For the classification task, we increased the state-of-the art by 15\% in accuracy.

In addition to BERT, we also implemented other pre-training language models like ELMo and ULMFit for comparison purposes. ULMFit, further pre-trained on a financial corpus, beat the previous state-of-the art for the classification task, only to a smaller degree than BERT. These results show the effectiveness of pre-trained language models for a down-stream task such as sentiment analysis especially with a small labeled dataset. The complete dataset included more than 3000 examples, but FinBERT was able to surpass the previous state-of-the art even with a training set as small as 500 examples. This is an important result, since deep learning techniques for NLP have been traditionally labeled as too "data-hungry", which is apparently no longer the case.

We conducted extensive experiments with BERT, investigating the effects of further pre-training and several training strategies. We couldn't conclude that further pre-training on a domain-specific corpus was significantly better than not doing so for our case. Our theory is that BERT already performs good enough with our dataset that there is not much room for improvement that further pre-training can provide. We also found that learning rate regimes that fine-tune the higher layers more aggressively than the lower ones perform better and are more effective in preventing catastrophic forgetting. Another conclusion from our experiments was that, comparable performance can be achieved with much less training time by fine-tuning only the last 2 layers of BERT. 

Financial sentiment analysis is not a goal on its own, it is as useful as it can support financial decisions. One way that our work might be extended, could be using FinBERT directly with stock market return data (both in terms of directionality and volatility) on financial news. FinBERT is good enough for extracting explicit sentiments, but modeling implicit information that is not necessarily apparent even to those who are writing the text should be a challenging task. Another possible extension can be using FinBERT for other natural language processing tasks such as named entity recognition or question answering in financial domain. 

\section{Acknowledgements}
I would like to show my gratitude to Pengjie Ren and Zulkuf Genc for their excellent supervision. They provided me with both independence in setting my own course for the research and valuable suggestions when I need them. I would also like to thank Naspers AI team, for entrusting me with this project and always encouraging me to share my work. I am grateful to NIST, for sharing Reuters TRC-2 corpus with me and to Malo et al. for making the excellent Financial PhraseBank publicly available.

\end{document}

